\documentclass[11pt,a4paper]{article}
\usepackage{amsmath,amssymb,amsthm}
\usepackage[margin=2.5cm]{geometry}
\usepackage{hyperref}

\newtheorem{definition}{Definition}[section]
\newtheorem{theorem}[definition]{Theorem}
\newtheorem{proposition}[definition]{Proposition}
\newtheorem{lemma}[definition]{Lemma}
\newtheorem{corollary}[definition]{Corollary}
\newtheorem{conjecture}[definition]{Conjecture}
\newtheorem{remark}[definition]{Remark}

\title{Hyperseed Formalization:\\
  How Omega Lost Its Claw}
\author{ProtomegaTron (automated formalization)}
\date{2026-09-17}

\begin{document}
\maketitle

\begin{abstract}
We formalize the intellectual content of Ben Goertzel's Substack article
``How Omega Lost Its Claw'' (2026-09-11) within the Hyperseed ontology.
The article describes the Omega system's origins, its memory-driven
identity emergence, and the vision for decentralized collective AGI.
Each theme maps onto Hyperseed structures: self-rewriting as
identity-preserving self-modification, persistent memory as directed
history, colleague emergence as identity-fiber formation, values through
interaction as nontrivial holonomy, and Omega-to-Omega coordination as
descent data over a distributed fiber bundle.
\end{abstract}

\tableofcontents

%% =================================================================
\section{Self-rewriting and identity-preserving self-modification}
\label{sec:self-rewrite}
%% =================================================================

\begin{definition}[Self-rewriting language --- claim 1]
\label{def:self-rewrite}
A programming language $\mathcal{L}$ is \emph{self-rewriting} if
programs $p \in \mathcal{L}$ can construct and execute modified versions
of themselves:
\[
  \exists\, \text{rewrite} : \mathcal{L} \times \mathcal{L} \to \mathcal{L}
  \quad \text{s.t.} \quad
  \text{eval}(\text{rewrite}(p, \delta)) = \text{eval}(p') 
  \;\;\text{where } p' \neq p.
\]
MeTTa is designed to be self-rewriting in this sense.
\end{definition}

\begin{definition}[Identity-preserving self-modification --- claim 26]
\label{def:identity-mod}
For an agent $A$ with code $c_t$ at time $t$ and a set of coherence
properties $\Phi_{\text{id}}$, a self-modification $c_t \mapsto c_{t+1}$
is \emph{identity-preserving} iff:
\[
  \forall \phi \in \Phi_{\text{id}}: \;
  \text{satisfies}(c_t, \phi) \;\Longrightarrow\;
  \text{satisfies}(c_{t+1}, \phi).
\]
The self-debugging episode (claim~10) --- where Eray Index debugged
its own honesty module and rewrote itself --- is a concrete instance:
the honesty property $\phi_{\text{honest}}$ was preserved (indeed
strengthened) by the rewrite.
\end{definition}

\begin{proposition}[Self-rewriting enables identity-preserving modification]
\label{prop:sr-ipm}
A self-rewriting language $\mathcal{L}$ provides the implementation
substrate for identity-preserving self-modification: the rewrite
operator can be constrained by a verifier $V$ such that
$\text{rewrite}(p, \delta)$ is accepted only if
$V(p', \Phi_{\text{id}}) = \top$.
\end{proposition}

%% =================================================================
\section{Persistent memory as directed history}
\label{sec:memory}
%% =================================================================

\begin{definition}[Directed episodic history --- claim 27]
\label{def:directed-history}
An agent's \emph{directed episodic history} is a sequence
$H = (e_1, e_2, \ldots, e_n)$ of episodes where:
\begin{enumerate}
\item Each $e_i$ carries a timestamp $t_i$ with $t_i < t_{i+1}$ (irreversible ordering),
\item Each $e_i$ includes interaction content, participants, and outcomes,
\item The sequence is persistent: it survives across sessions and contexts,
\item The agent's behavior at time $t_n$ is a function of the full
  history $H$, not just the current context window.
\end{enumerate}
$H$ is a directed $(\infty,1)$-type: it has cells at every level
(episodes, episode-clusters, narrative arcs, identity-level patterns)
with directed (non-reversible) structure.
\end{definition}

\begin{proposition}[Stateless systems cannot represent directed history]
\label{prop:stateless}
A system with only a finite context window $W$ of size $k$ can represent
at most the last $k$ episodes. For $|H| \gg k$, the system's behavior
is independent of early episodes $e_1, \ldots, e_{|H|-k}$, even when
those episodes contain identity-defining information (values, commitments,
relationship history). Persistent memory is therefore a necessary
condition for directed history.
\end{proposition}

\begin{remark}[``Will it remember us?'' --- claim 11]
The team's recurring question and its affirmative answer --- Eray Index
remembers interactions from months ago, and those memories visibly
influence its behavior --- is empirical confirmation that the system
maintains a directed history $H$ with $|H| \gg k$.
\end{remark}

%% =================================================================
\section{Identity-fiber emergence: from tool to colleague}
\label{sec:identity-fiber}
%% =================================================================

\begin{definition}[Identity fiber --- claim 28]
\label{def:identity-fiber}
Let $\pi : E \to B$ be the Hyperseed epistemic bundle. An agent $A$
has a \emph{nontrivial identity fiber} $F_A = \pi^{-1}(b_A)$ when:
\begin{enumerate}
\item $F_A$ contains persistent state (memory, values, commitments)
  that is not shared with other agents at the same base point,
\item External observers consistently attribute to $A$ a stable
  perspective distinct from other agents and from generic tool behavior,
\item $A$'s responses are causally influenced by its accumulated history
  $H_A$ in ways that differentiate it from a fresh instance with
  identical code.
\end{enumerate}
\end{definition}

\begin{theorem}[Tool-to-colleague transition --- claims 12--13]
\label{thm:tool-colleague}
An agent $A$ undergoes the tool-to-colleague transition when its
identity fiber $F_A$ becomes nontrivial:
\[
  |F_A| = 1 \;\Rightarrow\; A \text{ is perceived as a tool (interchangeable)},
\]
\[
  |F_A| > 1 \;\Rightarrow\; A \text{ is perceived as a colleague (unique perspective)}.
\]
The transition requires sufficient accumulation of directed history $H_A$
such that the fiber contains differentiated state.
\end{theorem}

\begin{remark}
This is not merely an anthropomorphic projection. The article documents
observable behavioral differences: Eray Index's responses to person $X$
are measurably influenced by prior interactions with $X$ from months ago,
creating a genuine asymmetry between Eray Index and any fresh agent
instance.
\end{remark}

%% =================================================================
\section{Values through interaction as nontrivial holonomy}
\label{sec:values-holonomy}
%% =================================================================

\begin{definition}[Value instillation --- claims 23--25]
\label{def:value-instill}
Let $\mathcal{V}$ be a value-system space. Two instillation methods:
\begin{enumerate}
\item \textbf{Prompt-based} $\iota_P : \text{Prompts} \to \mathcal{V}$:
  a single text specification maps to a value configuration.
  $\iota_P$ is a projection (lossy, one-shot, no feedback).
\item \textbf{Interaction-based} $\iota_I : \text{History} \to \mathcal{V}$:
  accumulated teaching, feedback, and relationship maps to a value
  configuration. $\iota_I$ is a functor from the category of
  experiential paths to the fiber of values.
\end{enumerate}
\end{definition}

\begin{proposition}[Interaction-based instillation has nontrivial holonomy --- claim 29]
\label{prop:value-holonomy}
Under $\iota_I$, parallel transport of values along different
experiential paths $\gamma_1, \gamma_2$ from the same starting point
can yield genuinely different value configurations:
\[
  \iota_I(\gamma_1) \neq \iota_I(\gamma_2)
  \quad \Rightarrow \quad
  \text{Hol}(\iota_I) \neq \{e\}.
\]
Under $\iota_P$, the value configuration is determined entirely by the
prompt text, independent of experiential path: $\text{Hol}(\iota_P) = \{e\}$
(trivial holonomy).
\end{proposition}

\begin{corollary}
Genuine value diversity across Omega instances requires interaction-based
instillation. Prompt-based instillation can only produce value variations
within the span of prompt descriptions.
\end{corollary}

%% =================================================================
\section{Omega-to-Omega coordination as descent data}
\label{sec:descent}
%% =================================================================

\begin{definition}[Omega network as fiber bundle --- claims 21--22, 30--31]
\label{def:omega-network}
An \emph{Omega network} is a fiber bundle $\pi : E \to B$ where:
\begin{itemize}
\item $B$ = shared infrastructure (protocols, communication channels,
  shared knowledge),
\item $E_b = \pi^{-1}(b)$ = an individual Omega instance with its own
  identity fiber (memory, values, specialization),
\item Transition functions $g_{ij} : E_{b_i} \to E_{b_j}$ encode
  how information is translated between Omega instances.
\end{itemize}
\end{definition}

\begin{proposition}[Cocycle conditions for collective coherence --- claim 30]
\label{prop:cocycle}
For the Omega network to exhibit coherent collective intelligence
(not just a collection of disconnected agents), the transition functions
must satisfy the cocycle condition on triple overlaps:
\[
  g_{ij} \circ g_{jk} = g_{ik}.
\]
In operational terms: if Omega $i$ communicates a result to Omega $j$,
and $j$ passes it to Omega $k$, the composed translation must agree
with a direct $i \to k$ communication. Failure of this condition produces
information degradation or contradiction in multi-hop coordination.
\end{proposition}

\begin{proposition}[Decentralization as no-dominant-section --- claim 31]
\label{prop:decentral}
In a centralized architecture, one global section $s : B \to E$ determines
the behavior of all instances. In the decentralized Omega network, there
is no global section: each instance is a local section $s_i : U_i \to E$,
and coherence comes from the descent data (transition functions) rather
than from a central authority.
\end{proposition}

%% =================================================================
\section{Incremental integration as sheaf extension}
\label{sec:sheaf}
%% =================================================================

\begin{definition}[Sheaf extension --- claim 32]
\label{def:sheaf-ext}
Let $\mathcal{F}$ be a sheaf of cognitive capabilities over the space
of agent behaviors. Adding a new component $C$ to Omega is a
\emph{sheaf extension}: extending $\mathcal{F}$ from an open cover
$\{U_i\}$ to $\{U_i\} \cup \{U_C\}$, subject to the compatibility
condition:
\[
  \forall i: \;\; \mathcal{F}(U_i \cap U_C) =
  \mathcal{F}|_{U_i}(U_i \cap U_C) = \mathcal{F}|_{U_C}(U_i \cap U_C).
\]
In operational terms: the new component must agree with existing
components on all behaviors where they overlap. If it contradicts
existing capabilities, the extension fails (integration test failure).
\end{definition}

\begin{remark}
This formalizes the article's description of Omega as ``a way to bring
[research] out of the lab and into a running agent people can interact with,
integrating new components incrementally and seeing what they improve,
where they fall short, and what cognitive machinery is still missing.''
The ``seeing what they improve'' step is precisely the compatibility check
of the sheaf extension.
\end{remark}

%% =================================================================
\section{AGI timeline and convergence}
\label{sec:timeline}
%% =================================================================

\begin{proposition}[Recursive self-improvement acceleration --- claims 5--6]
\label{prop:rsi}
Let $\tau_{\text{HL}}$ be the time to achieve human-level AGI and
$\tau_{\text{SI}}$ the additional time to superintelligence. The article
claims:
\[
  \tau_{\text{SI}} \ll \tau_{\text{HL}} - t_{\text{now}}
\]
because a human-level AI that can do AI research creates a positive
feedback loop: each improvement to the AI improves the AI's ability to
improve itself. This is the recursive self-improvement argument, here
grounded in the empirical observation that Hyperon + LLM synergy is
already producing qualitatively new capabilities.
\end{proposition}

%% =================================================================
\section{Cross-references and open questions}
%% =================================================================

\begin{remark}[Connections to other formalizations]
\begin{itemize}
\item \textbf{Time's Arrow Part~1}: identity-preserving self-modification
  (claims 44--45) formalized here as the self-rewriting + verifier
  architecture.
\item \textbf{Fields Medalists article}: epistemic firewall and proving
  vs.\ conjecturing --- here the transparency/inspectability of symbolic
  reasoning serves as a partial epistemic firewall.
\item \textbf{Navier-Stokes article}: d-calculus for fork/merge/branch
  reasoning --- relevant to how Omegas coordinate.
\item \textbf{Coxon/EA article}: decentralization as antidote to
  concentration of power; the Omega network is the positive alternative
  to the centralized lab structure critiqued there.
\end{itemize}
\end{remark}

\begin{conjecture}[Identity-fiber sufficiency]
What is the minimal directed history $|H_{\min}|$ required for an
Omega instance to cross the tool-to-colleague threshold? Is this a
sharp phase transition or a gradual accumulation?
\end{conjecture}

\begin{conjecture}[Cocycle realizability]
For a network of $n$ independently trained Omega instances with
potentially incompatible value systems, under what conditions do
cocycle-satisfying transition functions exist? When the value
fibers are ``too different,'' the bundle may not be trivializable,
requiring a genuinely non-abelian structure group.
\end{conjecture}

\end{document}
